\section{Method}
\label{sec:method}

\subsection{Architecture}
\label{sec:method:arch}

EEG-LeJEPA is a Joint-Embedding Predictive Architecture composed of three
components: a per-patch encoder, a causal Transformer predictor, and a
distribution-matching regularizer (SIGReg). The total parameter count is
2.86\,M, of which 1.08\,M live in the encoder and 1.78\,M in the predictor.

\paragraph{Patch tokenizer.}
Each 4\,s EEG window at 200\,Hz is a tensor of shape
$x \in \mathbb{R}^{C \times T}$ with $C{=}64$ channels and $T{=}800$ samples. A
1-D convolution with kernel and stride both equal to the patch size
$p{=}40$ samples (200\,ms) jointly mixes channels and time within each patch,
yielding a sequence of $T/p{=}20$ patch tokens, each of dimension $d{=}192$.
A sinusoidal positional embedding is added to the token sequence.

\paragraph{Per-patch encoder.}
Each token then passes through a two-layer per-patch MLP with GELU
activations, layer-normalized at the input of each MLP layer for stability.
The crucial design constraint is that the encoder \emph{must not} mix
information across patches: if it did, the encoder embedding $z_t$ would
leak information about $z_{t+1}$ and the LeJEPA next-token prediction
objective would degenerate into a trivial identity. We enforce this by
using only \texttt{nn.Linear} and pointwise nonlinearities in the
post-tokenizer encoder layers; there is no self-attention in the encoder.

\paragraph{Final encoder normalization: BatchNorm, not LayerNorm.}
The encoder closes with a BatchNorm1d projector. This choice is load-bearing
for SIGReg. LayerNorm forces per-sample unit norm, which fixes each
embedding's L2-magnitude to a constant and destroys the distribution-matching
signal that SIGReg requires; BatchNorm preserves per-sample magnitude
variation across the batch while keeping the per-coordinate statistics
well-conditioned~\cite[App.~A]{lewm2026}.

\paragraph{Causal Transformer predictor.}
The predictor consumes the encoder's token sequence and predicts the
next-step embedding. It is a 4-layer causal Transformer with embedding
dimension 192, four attention heads, MLP ratio 4, and dropout 0.1. Position
$t$ of its output is trained to match the encoder's embedding at position
$t{+}1$ under mean-squared error:
\begin{equation}
  L_\text{pred} = \frac{1}{N(T-1)}\sum_{i,t}
    \|\text{predictor}(z_i)_t - z_{i,t+1}\|_2^2,
  \label{eq:lpred}
\end{equation}
where $z_i \in \mathbb{R}^{T \times d}$ is the encoder output for sample $i$.
The predictor closes with the same BatchNorm1d normalization as the encoder.
Per LeJEPA~\cite{lejepa2025} and LeWorldModel~\cite{lewm2026}, there is
\emph{no} stop-gradient, \emph{no} EMA target encoder, and \emph{no}
teacher--student split: the same encoder produces both inputs and targets
for the prediction loss, and SIGReg alone prevents the trivial-solution
collapse that this would otherwise cause.

\subsection{SIGReg objective}
\label{sec:method:sigreg}

The training loss is
\begin{equation}
  L = L_\text{pred} + \lambda \cdot L_\text{SIGReg},
  \label{eq:ltotal}
\end{equation}
where $L_\text{SIGReg}$ is the Sketched Isotropic Gaussian Regularizer
of~\cite{lejepa2025}. Given the flattened embedding matrix
$Z \in \mathbb{R}^{NT \times d}$, we sample $M$ unit vectors
$u^{(m)} \sim \mathcal{U}(S^{d-1})$ uniformly on the unit sphere and apply a
univariate Gaussianity test $T(\cdot)$ to each one-dimensional projection
$Z u^{(m)} \in \mathbb{R}^{NT}$:
\begin{equation}
  L_\text{SIGReg}(Z) = \frac{1}{M}\sum_{m=1}^{M} T\bigl(Z u^{(m)}\bigr).
  \label{eq:sigreg}
\end{equation}
By the Cram\'er--Wold theorem~\cite{cramerwold1936}, the joint distribution
of $Z$ is standard normal if and only if every one-dimensional projection
is standard normal, so driving $L_\text{SIGReg}$ to zero is equivalent to
enforcing joint isotropic Gaussianity. We re-sample the unit vectors at
every training step so the regularizer's stochasticity does not bias the
encoder toward any particular fixed sketch.

\paragraph{Univariate test: Cram\'er--von Mises vs.\ Epps--Pulley.}
LeJEPA recommends the Epps--Pulley characteristic-function statistic for
$T(\cdot)$. We instead use the Cram\'er--von Mises (CvM) statistic, which
has a closed-form expression for samples sorted in ascending order
$x_{(1)} \le \dots \le x_{(n)}$:
\begin{equation}
  W^2 = \frac{1}{12n}
        + \sum_{i=1}^{n}\!\left[\,F(x_{(i)}) - \frac{2i-1}{2n}\right]^2,
  \label{eq:cvm}
\end{equation}
where $F$ is the standard normal CDF. CvM is sort-based, numerically
well-conditioned, and easily verifiable against
\texttt{scipy.stats.cramervonmises}; in our preliminary calibration both
statistics produced equivalent optimisation trajectories and we adopt CvM
for its implementation simplicity. The complete SIGReg loop is summarised
in Algorithm~\ref{alg:sigreg}.

\begin{algorithm}[t]
\caption{SIGReg loss with CvM univariate statistic.}
\label{alg:sigreg}
\begin{algorithmic}[1]
\REQUIRE Embeddings $Z \in \mathbb{R}^{NT \times d}$, number of slices $M$.
\STATE Sample $u^{(1)}, \dots, u^{(m)} \sim \mathcal{U}(S^{d-1})$ i.i.d.
\STATE For each $m$, project $\xi^{(m)} \gets Z u^{(m)}$ (vector of length $NT$).
\STATE For each $m$, sort $\xi^{(m)}$ ascending and compute $W^2_m$ via
       Eq.~\eqref{eq:cvm}.
\STATE $L_\text{SIGReg} \gets \tfrac{1}{M}\sum_m W^2_m$.
\RETURN $L_\text{SIGReg}$.
\end{algorithmic}
\end{algorithm}

\paragraph{Hyperparameters.}
We use $\lambda = 1.0$ (we ablate it across $\{0, 0.1, 0.3, 1.0, 3.0,
10.0\}$ in \cref{sec:results:lambda}) and $M = 1024$ slices, both calibrated
on a 3-subject 200-step pilot before the main runs.

\subsection{Linear probe protocol}
\label{sec:method:probe}

We evaluate the pretrained encoder via the standard SSL frozen-feature
linear-probe protocol. For each EEG epoch, the encoder produces a sequence
of $T_p$ patch embeddings of dimension $d$. We consider three feature
sources:
\begin{description}
  \item[\texttt{encoder\_mean}] mean-pool the per-patch encoder embeddings
    along the token axis; no temporal context.
  \item[\texttt{predictor\_mean}] mean-pool the predictor's hidden states;
    each position has seen the causal context up to that step.
  \item[\texttt{both\_mean}] concatenate the two preceding feature vectors.
\end{description}
The resulting feature vector is fed into a \texttt{scikit-learn}
\texttt{LogisticRegression} with $C{=}1.0$ and a maximum of 2{,}000
iterations. We evaluate cross-subject by leave-one-subject-out (LOSO)
cross-validation across 20 EEGMMIDB subjects (subjects 1--20). We report
the mean and standard deviation of per-fold accuracy, the macro one-vs-rest
AUC, and for multi-class tasks balanced accuracy and macro-$F_1$. A
randomly initialized identical-architecture model serves as the
no-pretraining control.

\subsection{Datasets}
\label{sec:method:data}

\paragraph{EEGMMIDB} (PhysioNet motor-imagery database; Schalk et
al.~\cite{eegmmidb2004}). 109 subjects, 64 EEG channels in the 10--10
montage, native sampling rate 160\,Hz, motor-imagery runs 4, 8 and 12
(left-hand vs.\ right-hand fist imagery). We resample to 200\,Hz, apply a
1--70\,Hz band-pass filter (with the upper cutoff clamped below the
source Nyquist for the few 128\,Hz subjects), a 60\,Hz notch filter, an
average reference, and finally epoch the recordings into 4-second
non-overlapping windows for pretraining or 4-second event-locked windows
starting at each cue for downstream classification.

We pretrain on the EEGMMIDB MI runs (4, 8, 12) of 50 subjects (subjects
51--100) by default; the 20 evaluation subjects (subjects 1--20) are
strictly held out from pretraining. We additionally report
ablations on 3, 20 and 109 pretraining subjects in
\cref{sec:results:scaling}. We probe two downstream tasks on
EEGMMIDB:
\emph{rest-vs-activity} (T0 cues against pooled T1+T2 cues, binary) and
\emph{left-vs-right MI} (T1 against T2, binary). Rest-vs-activity is the
easier sanity task; left-vs-right MI is the canonical BCI task and the
harder of the two.

\paragraph{BCI Competition IV Dataset 2a} (Brunner et
al.~\cite{bciiv2a}). 9 subjects, 22 EEG channels plus 3 EOG channels at
250\,Hz, 4-class motor imagery (left hand, right hand, both feet, tongue).
We use the training session (T) only. The pre-processing pipeline is
identical to EEGMMIDB except for a 50\,Hz notch (European mains) and
resampling to 200\,Hz. The 22 EEG channels are a subset of EEGMMIDB's 64,
which enables both naive channel-padded cross-dataset transfer (zero-pad
the 42 missing channels at inference) and channel-matched re-pretraining
(re-pretrain on the 22-channel intersection of EEGMMIDB). We probe the
full 4-class task by LOSO across the 9 subjects.
